\documentclass[11pt]{article}
\usepackage[margin=1in]{geometry}
\usepackage{amsmath,booktabs,graphicx,hyperref}
\title{Straight, Draw, and Fade: Shot Pattern Variance}
\author{UpstreamDrift}
\date{October 8, 2026}
\begin{document}
\maketitle
\section{Status and Scope}
The original four driver result bundles are historical controls from a
superseded impact approximation: tangential impulse changed spin but not
linear velocity. Corrected three-club experiments must replace their
conclusions. Fixed loft also excluded face--loft covariance.

This is an editable standalone experiment reference, separate from the governed
engineering design manual. The canonical manual remains the QMD source under
\texttt{manuals/upstreamdrift}. This experiment does not qualify the underlying
impact or flight model against measured shot populations. It estimates
conditional outcome variance from face-angle error, not a golfer's ability to
produce a chosen pattern. All results are first ground contact (carry); bounce
and roll are excluded.

\section{Inputs, Frames, and Pairing}
User-facing face and path angles are measured from the target, positive right
for a right-handed golfer. Straight uses $(f,p)=(0,0)$ degrees, draw uses
$(1.5,3)$, and fade uses $(-1.5,-3)$. The doubled-curve trials
use $(3,6)$ and $(-3,-6)$ while retaining the straight baseline. The draw face is therefore $1.5$ degrees
left of its path. The repository solver uses $x$ forward, $y$ left, and $z$ up;
the adapter changes the sign of horizontal input and output angles once.
The CSV diagnostic \texttt{spin\_axis\_tilt\_deg} instead retains the native
axis tilt convention, positive toward draw/left curvature. It is not
Trackman's right-positive spin-axis convention.
Every pattern receives the same 10,000 independent errors
$\epsilon_i\sim\mathcal{N}(0,\sigma_F^2)$, with separate trials at
$\sigma_F=1$ and $2$ degrees. In each trial,
$f_i=\bar f+\epsilon_i$ and $p_i=\bar p$. Shared errors provide paired
comparisons; finite draw and fade samples need not be exact mirrors.
The random seed is 20261008, using NumPy's default generator.

The corrected study uses three illustrative hybrid delivery presets:
\begin{center}
\begin{tabular}{lrrrrr}
\toprule
Club & Speed (m/s) & Loft ($^\circ$) & Attack ($^\circ$) & Lie ($^\circ$) & Mass (kg)\\
\midrule
Driver &45&12.8&-0.9&58.5&0.200\\
7-Iron &36&24.0&-4.0&63.0&0.272\\
Pitching Wedge &32&36.7&-5.0&64.0&0.300\\
\bottomrule
\end{tabular}
\end{center}
These are assumptions for comparative modeling, not measured player inputs or
reproduced specifications of one complete club. Trackman's dynamic-loft page
reports driver $12.8^\circ$ as a Tour average; its PW $36.7^\circ$ is an
optimizer example at 72 mph, not a measured Tour PW mean. The 7-iron
$24^\circ$ delivered loft is an explicit assumption. Manufacturer static lie
angles supply assumed shaft elevations. Speeds, attack angles, and masses
are illustrative inputs unless the source receipt separately identifies a
measurement. Nominal lean is zero; axis sensitivity explores 0, 10, and
20 degrees where the geometry is valid.

The experiment matrix is three clubs, two face standard deviations, two curve
magnitudes, and two delivery modes (fixed loft and assumed shaft rotation):
24 bundles, each with 10,000 shots per shape. The ball begins stationary
without spin; strike is central. Restitution, friction, ball mass and radius,
and atmospheric parameters are recorded in receipts. There is no wind,
speed error, strike error, or path error. Driver scoring uses a hypothetical
tee/fairway/rough context; 7-iron and PW scoring use approach/green contexts.
The earlier $45\,\mathrm{m/s}$, $10.9^\circ$ loft, zero-attack driver case is
retained only in superseded historical bundles.

\section{Corrected Central-Contact Impact Equations}
The experiment calls the public \texttt{ImpactSolverAPI} rigid-body model.
For unit face normal $\mathbf n$, relative centre velocity
$\mathbf v_r=\mathbf v_c^- -\mathbf v_b^-$, ball mass $m_b$, club mass $m_c$,
restitution $e$, ball radius $R$, and inertia $I_b=2m_bR^2/5$:
\[
 m_*=(m_b^{-1}+m_c^{-1})^{-1},\qquad
 J_n=(1+e)m_*\mathbf v_r\cdot\mathbf n.
\]
The ball contact point is $-R\mathbf n$ from its centre. Tangential surface
slip, including initial ball spin, and its effective inverse mass are
\[
 \mathbf u_t=\mathbf v_r-(\mathbf v_r\cdot\mathbf n)\mathbf n
                 +R(\boldsymbol\omega_b^-\times\mathbf n),\qquad
 D=m_b^{-1}+m_c^{-1}+R^2/I_b.
\]
For positive normal approach, Coulomb friction supplies
\[
 \mathbf J_t=\min\{\mu J_n,\|\mathbf u_t\|/D\}
                 \frac{\mathbf u_t}{\|\mathbf u_t\|},\qquad
 \mathbf J=J_n\mathbf n+\mathbf J_t.
\]
The tangential impulse is zero at zero slip or zero friction. Both linear
velocities and ball spin receive their corresponding impulse:
\[
 \mathbf v_b^+=\mathbf v_b^-+\mathbf J/m_b,\quad
 \mathbf v_c^+=\mathbf v_c^- -\mathbf J/m_c,\quad
 \boldsymbol\omega_b^+=\boldsymbol\omega_b^-+
                 (R/I_b)(\mathbf J_t\times\mathbf n).
\]
Earlier result bundles omitted $\mathbf J_t$ from both translational updates
and used the infinite-club sticking cap $(2/7)m_b\|\mathbf u_t\|$.
Those bundles are superseded historical controls. Regression tests first
failed against those defects, then passed with the correction. Independent
arbitrary-3D tests check linear and angular momentum, normal restitution,
nonincreasing kinetic energy, sticking contact, initial spin, and rotation
covariance. These establish internal consistency of central contact with a
translating club; off-centre club rotation, shaft response, and measured launch
accuracy remain unqualified.

\section{Geometric Face--Loft Coupling and Shaft Lean}
This is an explicitly assumed rigid geometry, not observed player data.
The native right-handed frame is $+X$ forward, $+Y$ left, and $+Z$ up.
User face angle is positive right. At zero twist, the delivered face normal is
$\mathbf n_0=(\cos L,0,\sin L)$ for nominal delivered loft $L$.
Given assumed shaft elevation $\ell$ and forward lean $a$:
\[
 s_x=\sin\ell\tan a,\qquad
 \mathbf s=(s_x,\sqrt{\cos^2\ell-s_x^2},\sin\ell).
\]
This construction is valid only when its square-root argument is nonnegative.
Rodrigues rotation about the unit shaft axis gives
\[
 \mathbf n(\theta)=\mathbf n_0\cos\theta+
 (\mathbf s\times\mathbf n_0)\sin\theta+
 \mathbf s(\mathbf s\cdot\mathbf n_0)(1-\cos\theta).
\]
Delivered face and loft are extracted as
\[
 F=\operatorname{atan2}(-n_y,n_x),\qquad
 L_d=\operatorname{atan2}(n_z,\sqrt{n_x^2+n_y^2}).
\]
Near zero twist,
\[
 \frac{dF}{d\theta}=-(s_z-s_x\tan L),\quad
 \frac{dL_d}{d\theta}=-s_y,\quad
 \frac{dL_d}{dF}=\frac{s_y}{s_z-s_x\tan L}.
\]
With zero lean, $dL_d/dF=\cot\ell$: closing tends to reduce loft in this
specific rotation mode. Rotation about the vertical axis instead preserves
loft. Face angle alone therefore does not identify delivered loft.

The main comparison applies this geometry to the face error $F-F_0$ around
each pattern's nominal face $F_0$, then rotates the axis and normal together
by the nominal target yaw. All shapes share the same zero-error delivered
loft, speed, attack angle, and strike. This isolates sensitivity to the assumed
face--loft coupling without giving one shape a different baseline loft.
Fixed-loft trials provide the corresponding control. Separate shaft-lean
sensitivity holds nominal loft fixed and changes only the axis; it does not
represent pitching a fixed physical club forward. That latter operation
changes nominal loft as well and must be analyzed separately.

Less delivered loft does not necessarily increase carry. Launch, spin, speed,
and aerodynamic trajectory jointly determine carry, so long-left/short-right
behavior is measured from the modeled endpoint covariance rather than imposed.
Static manufacturer lie is used only as an assumed delivered shaft elevation;
dynamic lie, shaft bending, and golfer covariance are unmeasured.

\section{Flight Integration and Landing}
Each shot calls the existing \texttt{BallFlightSimulator}, backed by the
\texttt{upstream-physics} Rust kernel. Launch velocity uses metres per second,
launch and azimuth angles radians, and spin magnitude RPM at the Python
boundary, converted to radians per second inside the kernel. The spin axis
is normalized. The solver integrates position, velocity and decaying spin
with RK4, drag and spin-induced aerodynamic force:
\[
 \dot{\mathbf r}=\mathbf v,\qquad
 m_b\dot{\mathbf v}=m_b\mathbf g+\mathbf F_D+\mathbf F_M,\qquad
 \dot\omega=-\lambda\omega.
\]
The production step is $0.02\,\mathrm{s}$ with a $12\,\mathrm{s}$ horizon.
Landing is linearly interpolated between the last above-ground and first
nonpositive-height samples. Failure to land raises an error; it never
silently treats the final airborne sample as a landing.
The Rust coefficient path is distinct from the constant-spin Waterloo/Penner
model in \texttt{flight\_models.py}; their calibrations must not be mixed.
The kernel applies a single spin force (not separate additive backspin lift
and sidespin Magnus forces). With $A=\pi R^2$ and relative airspeed $V$:
\[
 \mathbf F_D=-\tfrac12\rho C_D A V\mathbf v,\qquad
 S=R|\boldsymbol\omega|/V,\qquad
 C_L=\min(0.26,\;0.70 S^{0.645}),
\]
\[
 \mathbf F_M=\tfrac12\rho C_L A V^2
 \frac{\boldsymbol\omega\times\mathbf v}
 {|\boldsymbol\omega\times\mathbf v|}.
\]
The zero-cross-product case returns zero spin force. The shared drag-crisis
curve evaluates Reynolds number $\rho V(2R)/\mu$ with
$\mu=1.81\times10^{-5}\,\mathrm{Pa\,s}$ and base drag coefficient $0.25$,
using the exact source-bound constants in \texttt{aerodynamics.rs}.
Spin decay uses $\lambda=0.05\,\mathrm{s}^{-1}$.

\section{Aiming and Outcome Statistics}
Raw outcomes keep the requested target-relative angles. A second comparison
rotates each entire pattern by a single angle computed from its nominal,
zero-error landing, placing that nominal landing on the target line. This is
an aiming adjustment, not per-shot correction or a sample-mean fit. The
fixed target distance is the straight nominal carry; differences in curved
carry remain in the error metric.
\[
 \alpha=\operatorname{atan2}(y_0,x_0),\qquad
 \begin{pmatrix}x'\\y'\end{pmatrix}=
 \begin{pmatrix}\cos\alpha&\sin\alpha\\-\sin\alpha&\cos\alpha\end{pmatrix}
 \begin{pmatrix}x\\y\end{pmatrix}.
\]
Lateral variance is the unbiased sample variance with denominator $n-1$;
standard deviation is its square root. Central 90\% width is the empirical
95th minus 5th percentile. Target RMSE is
$\sqrt{n^{-1}\sum_i[(x_i-x_T)^2+y_i^2]}$. Target-hit probability counts
landings within a $15\,\mathrm m$ radius circle, not a fairway strip.
Paired percentile-bootstrap variance-ratio intervals resample common shot
indices 500 times. These intervals describe Monte Carlo error only; they
exclude model discrepancy and parameter uncertainty.

\section{Distance Control}
Carry is not fixed: face error changes impact and spin, and curved patterns
have shorter mean carry. The supplementary common-range diagnostic maps each
nominal-aimed endpoint to
\[
 (x'',y'')=\frac{x_T}{\sqrt{x'^2+y'^2}}(x',y').
\]
It preserves landing bearing and sets every radial range to the same straight
nominal distance. It is geometric normalization, not a new flight simulation,
and cannot isolate a causal physical mechanism. At one degree face SD, the
curved SD reduction falls from about 0.95\% at actual carry to 0.76\% at common
range; therefore distance explains part, but not all, of this modeled effect.

\section{Approach Scoring Estimate}
The scoring scenario is an approach from fairway at the fixed target distance
(about $195.334\,\mathrm m$), with the hole at the target and a circular green
of radius $15\,\mathrm m$. Nominal-aimed carry endpoints are treated as final
resting positions. Inside the green, the remaining-strokes benchmark depends
on putting distance; outside, it uses unobstructed rough distance. There are
no bunkers, water, out-of-bounds, slopes, bounce, or roll. Alternative green
radii expose sensitivity to this synthetic layout.
For historical-tour expected strokes $J(d,\ell)$ at distance $d$ and lie
$\ell$, the existing source-backed Tools analysis computes
\[
 SG_i=J(x_T,\mathrm{fairway})-1-J(d_i,\ell_i).
\]
All patterns share the starting state, so their relative benefit depends only
on the difference in expected strokes remaining. Reported benefits are per
approach, not per round. Multiplying by an assumed number of identical
approaches is an illustration, not a full-course simulation.
The historical fairway and rough values come from Broadie's published
benchmark table. The putting approximation and its source reconciliation are
recorded in the scoring provenance. In the 2011 preprint, literal evaluation
of the displayed three-putt formula with the listed coefficient order gives
negative probabilities and conflicts with the stated anchors; that literal
experiment was rejected. A reconciled approximation must not be represented
as an independently verified exact reproduction of the author's fit.
The implemented approximate three-putt expression uses putt distance $d_f$
in feet:
\[
 \widetilde p_3=\frac{1}{1+\exp(5.49-0.106d_f+0.000563d_f^2)}-0.00398,
 \qquad p_3=\min\{\widetilde p_3,1-p_1\}.
\]
The upper bound prevents a negative two-putt probability for tap-ins; the
published one-putt expression uses yards. The baseline supports only its
explicit distance range and rejects unsupported queries rather than extrapolating.
Sampling intervals exclude putting-model uncertainty, green-layout error,
and impact-model discrepancy.

\section{Interpretation and Observability Limits}
The relevant question is whether curved patterns reduce spread or target
error after equal aiming treatment, not whether their nominal balls curve
back toward the line. A symmetric, still-air model cannot establish a
physical draw-versus-fade asymmetry. A nonzero face-to-path mean also does
not eliminate opposite curvature: a $1.5$-standard-deviation offset leaves
about $6.68\%$ probability of crossing zero face-to-path under the assumed
normal error distribution. At two degrees SD, the theoretical crossing
probability rises to about $22.66\%$. Real golfers may have correlated face/path
errors, different variability by intention, strike-dependent gear effect,
wind, varying speed or loft, and asymmetric hazards. Those require measured
inputs and additional experiments.
With only one uncertain input, the landing set is a narrow one-dimensional
curve in the horizontal plane, not a measured two-dimensional dispersion
ellipse. No ellipse-area advantage is inferred from this experiment.
For intuition, a local linear landing model
$Y\simeq aF+b(F-P)$ with fixed path gives
$\operatorname{Var}(Y)\simeq(a+b)^2\sigma_F^2$, independent of the mean
face and path if the sensitivities are constant. Any simulated spread
difference therefore arises from changing nonlinear sensitivities, not the
mere existence of a nonzero curvature mean. This linearization explains the
comparison; it is not used to generate any trajectory or landing.

\section{Evidence and Reproduction}
The adjacent README and checked-in result bundle record numerical results,
validation commands, dependency/source hashes and graphic provenance.
The graphics show a limited set of representative flights while the
dispersion and statistics use all 10,000 shots per pattern. Full reproduction:
\begin{verbatim}
python3 -m src.tools.shot_pattern_analysis \
  --shots 10000 --seed 20261008 \
  --output docs/research/shot_pattern_analysis/results
python3 -m src.tools.shot_pattern_analysis \
  --shots 10000 --seed 20261008 --face-sd 2 --curve-scale 2 \
  --output docs/research/shot_pattern_analysis/results_large_curve_sd2
\end{verbatim}
The native Rust dependency must be installed from this checkout. Compilation
of this reference is a document check, not scientific acceptance.

\section{External Context}
Trackman defines face-to-path as the horizontal difference between face and
path and relates it to curvature; see
\url{https://www.trackman.com/blog/face-to-path} and
\url{https://www.trackman.com/blog/spin-axis}.
Historical strokes-gained methodology and benchmarks are described in
Broadie, \emph{Assessing Golfer Performance on the PGA TOUR},
\url{https://business.columbia.edu/sites/default/files-efs/pubfiles/4996/assessing_golfer_performance.full.pdf}.
These definitions support the input interpretation, not validation of this
repository's numerical predictions.
\end{document}
